\section{Introduction}\label{sec:introduction}

For $n\ge3$, a proper random 3-clause chooses three distinct variables
uniformly from $x_1,\ldots,x_n$ and gives each variable an independent
fair sign. Let $\Phi(n,m)$ be the conjunction of $m\ge0$ independent
clauses with this law. Thus its clauses are sampled with replacement
from the $8\binom n3$ possibilities; the formula with no clauses is
satisfiable. Write
\[
 p(n,m)=\Prob\{\Phi(n,m)\text{ is satisfiable}\}.
\]
The ratio $m/n$ is the clause density. We prove that the limiting density
separating satisfiable and unsatisfiable formulas can be approximated to
any prescribed accuracy by one finite algorithm.

Carenini~\cite[Corollary~1.2]{Carenini2026Polynomial} establishes the
existence of a limiting satisfiability threshold for every fixed
$k\ge3$. We credit Carenini with priority for the resolution of the
satisfiability conjecture. The contribution here is computability of
the 3-SAT threshold, beyond its existence.

\begin{theorem}\label{thm:main}
There is a real number $\alpha_3\in(0,\infty)$ such that, for every fixed
real $a\ge0$,
\[
 \lim_{n\to\infty}p(n,\lfloor an\rfloor)=
 \begin{cases}1,&a<\alpha_3,\\0,&a>\alpha_3.\end{cases}
\]
There is one finite deterministic Turing machine which, on input $1^r$
for any integer $r\ge1$, halts with a rational number $q_r$ satisfying
$|q_r-\alpha_3|\le2^{-r}$. Its only input is the unary precision $1^r$;
it uses no oracle, advice, or noncomputable real constant.
\end{theorem}

The strict-side limits determine $\alpha_3$ uniquely, since two different
constants would give opposite limits at a density between them. The
theorem leaves the behavior at $a=\alpha_3$ open. Its algorithm computes
the limiting density; it does not search for satisfying assignments to
an input formula. We obtain no efficiency bound, and the computability
argument concerns clause size three.

Random 3-SAT has long served as a test case for the typical behavior of
satisfiability problems. Experiments of Mitchell, Selman, and
Levesque~\cite{MSL1992} placed difficult instances for the Davis--Putnam
procedure near the density where half the formulas were satisfiable.
Algorithmic and counting arguments then constrained the transition's
location. Kaporis, Kirousis, and Lalas~\cite{KaporisKirousisLalas2006}
obtained the lower bound $3.52$, independently reported by Hajiaghayi
and Sorkin~\cite{HajiaghayiSorkin2003}; D\'iaz, Kirousis, Mitsche, and
P\'erez-Gim\'enez~\cite{DiazKirousisMitschePerezGimenez2009} obtained the
upper bound $4.4898$. These bounds bracket any limiting threshold.
From statistical physics, M\'ezard, Parisi, and Zecchina developed the
cavity and survey-propagation description~\cite{MPZ2002}. The
calculations of M\'ezard and Zecchina~\cite{MZ2002}, refined by Mertens,
M\'ezard, and Zecchina~\cite{MMZ2006}, predict a threshold near $4.267$.
Theorem~\ref{thm:main}
establishes existence and computability without evaluating the constant
or identifying it with that prediction.

Two further distinctions explain the mathematical task. A transition
window can become narrow while its center still varies with the system
size. Friedgut's theorem, with Bourgain's appendix~\cite{Friedgut1999},
establishes a sharp transition for every fixed $k\ge3$ about a
size-dependent location. In earlier work,
Carenini~\cite[Corollaries~7.10--7.11]{Carenini} obtained an
$O(n/\log n)$ central clause window for each fixed $k\ge3$, including
the model with independently sampled proper clauses. Her polynomial
improvement to $O_{k,\eta}(n^{1/2+1/k})$ between probability levels
$\eta$ and $1-\eta$, for each fixed $0<\eta<1/2$, in
\cite[Theorem~1.1]{Carenini2026Polynomial} yields limiting thresholds
for every fixed $k\ge3$.
Bayati, Gamarnik, and Tetali~\cite{BGT2013} developed sparse
interpolation and proved limits for normalized optimization values and
finite-temperature normalized log partition functions. Their near-satisfiability
conclusion permits $o(n)$ violated clauses; exact satisfiability requires
none. Ding, Sly, and Sun~\cite{DSS2022} proved the
limiting-threshold conjecture, including the predicted value, for all
sufficiently large fixed clause sizes. The probability argument below
also proves existence at clause size three as part of the development
of the computability result.

Even when a finite-size limit exists, it need not be a computable real.
Specker~\cite[Satz~IV]{Specker1949} constructed a bounded nondecreasing
computable sequence of rationals with a noncomputable limit. Our proof
supplies two families of finite certificates. The first certifies
rational lower bounds on $\alpha_3$ that approach it arbitrarily closely.
The second certifies rational upper bounds, with a certificate at every
rational density strictly above $\alpha_3$. A fair search through both
families produces certified intervals of arbitrarily small length. Its
stopping rule uses the certified interval itself and requires no
convergence modulus for the finite-size quantities.
